%==============================================================================
% Conclusion + Ethics/Broader-Impact + Reproducibility + LLM-use disclosure
%==============================================================================
% This file is an \input-able section. The conclusion proper is set as
% \section{Conclusion} so it can replace the lorem stub in main.tex; the three
% required closing blocks follow as unnumbered \section* blocks, matching the
% main.tex skeleton's placement (after the conclusion, before References).
%==============================================================================

\section{Conclusion}
\label{sec:conclusion}

Organizational memory is a distinct problem from the conversational and
personal-memory regimes that current benchmarks measure: it is graph-shaped,
bi-temporal, and decision-traced, and these are precisely the properties on
which the strongest conversational systems are no longer tested. \bench{}
isolates them over a synthetic corpus whose ground truth is
projected deterministically from a known source-of-truth graph rather than read
back out of free text.

Measured against this task, the field is far from a solution. Across the four
publicly available memory systems we evaluate (full numbers in
\S\ref{sec:experiments}, Table~\ref{tab:main-results}), no system clears a low
rubric-weighted ceiling, and every
system collapses on bi-temporal and contradiction questions---the categories that
most directly express what makes organizational memory hard. We further find that
the small-tier ranking does not survive the move to organizational scale: of the
three trailing systems, two swap rank between tiers, while the front-runner
extends its lead. Conversational-scale behavior, in other words, does not predict
organizational-scale behavior. The headline of this paper is therefore not that
one system wins. It is that organizational-scale memory is an \emph{unsolved}
problem with large remaining headroom, and that the gap is concentrated in
exactly the temporal and provenance-tracing abilities that long-horizon agentic
work will depend on.

We also contribute the evaluation methodology that makes these numbers credible.
Fair comparison of hosted, \emph{asynchronous} memory systems requires a
readiness gate: a naive ingest-returned or count-plateau signal can query an
effectively empty store and silently understate a system, so hosted systems must
be drained on a probe search to a stable non-zero state before scoring. We report
this fix, its per-system effect, and the residual variance it does not remove.

We release everything needed to build on this work: the \bench{} corpus and
question sets for both shipped tiers, the source-of-truth graphs and rubrics, the
uniform evaluation harness with each system's exact configuration, and a planned
public leaderboard. We invite the community to add systems, contest our
configurations, and push the medium and larger tiers. The immediate next steps
are clear: extend the corpus to additional organizations and to the larger token
tiers to trace the scaling curve, and treat closing the temporal and
justification-chain gap---not
incremental gains on saturated conversational benchmarks---as the open frontier
for agent memory.

%------------------------------------------------------------------------------
% ETHICS AND BROADER IMPACT  (required at NeurIPS; expected elsewhere)
%------------------------------------------------------------------------------
\section*{Ethics and Broader Impact}
\label{sec:ethics}

\bench{} contains no real organizational data. The seed organization (Helix
Logistics), its personnel, customers, and every artifact are entirely synthetic,
generated by a local open-weight model from configured scale parameters; no
real company's confidential communications, and no personally identifiable
information about real individuals, are present. This sidesteps the consent and
privacy concerns that attach to corpora scraped from genuine enterprise systems,
and the freshly generated fictional organization gives the corpus the
contamination resistance argued in \S\ref{sec:design:why-synthetic}.

The benchmark's broader purpose is to make organizational-memory failures
\emph{measurable} before agents are widely deployed on organizational context:
the failure we most want to surface is not ``the agent does not know'' but ``the
agent confidently acts on a superseded or contradicted fact,'' which our
per-category results show is common today. Two caveats bound the scope.
Optimizing directly to \bench{} risks overfitting to its generation pipeline and
rubric structure, and its single English-language organization may not generalize
across languages, industries, or org structures; we therefore release multiple
tiers, commit to refreshed and held-back instances, and read \bench{} as a
diagnostic of where systems fail rather than a target to maximize in isolation.
And because a system that retrieves organizational context well is dual-use---the
same capability could, if mis-scoped, surface information beyond a user's access
boundary---scope-aware retrieval is part of the task definition rather than an
afterthought.

%------------------------------------------------------------------------------
% REPRODUCIBILITY STATEMENT
%------------------------------------------------------------------------------
\section*{Reproducibility Statement}
\label{sec:reproducibility}

We release the full \bench{} corpus, question sets, source-of-truth graphs, and
weighted rubrics for both shipped tiers, together with the uniform evaluation
harness, which records each system's exact configuration so that every reported
number is reconstructable. Ground truth is projected deterministically from the
source-of-truth graph, so the answer key is reproducible from the released graphs
alone, independent of any model. The single-seed, small-$n$, LLM-judged scope of
the reported runs (and the resulting caution against over-reading sub-0.1
differences) is stated in full in \S\ref{sec:limitations}. Full per-system
hyperparameters, version pins, the judge and answerer configurations, and the
per-tier change logs with pre-patch audit snapshots are given in
Appendix~\ref{app:hyperparams}.

%------------------------------------------------------------------------------
% USE OF LARGE LANGUAGE MODELS  (disclose proactively)
%------------------------------------------------------------------------------
\section*{Use of Large Language Models}
\label{sec:llm-use}

We disclose LLM use proactively, as required by current venue policies, in three
distinct roles. First, the \bench{} corpus is \emph{entirely LLM-generated} by a
local open-weight model (\texttt{gpt-oss:20b}), with ground-truth answers
projected deterministically from a structured source-of-truth graph rather than
authored by the model. Second, the evaluation pipeline uses Claude Sonnet~4.6 as
the rubric-based judge (with extended thinking) and, for systems lacking a native
synthesizer, as the shared neutral answerer. Third, the authors used an LLM to
assist with editing and clarity. The authors take full responsibility for all
scientific claims, experimental design, results, and analysis, which were
produced and verified by the authors; no LLM is an author of this paper.
